% =============================================================================
% APPENDIX SU: LIT-SUTTER: Team Decision Making and Developmental Economics
% =============================================================================
% Module: BCM2_LIT_SUTTER
% Category: LIT
% Version: 1.0 (January 2026)
% Status: Final
%
% Comprehensive integration of Matthias Sutter's research contributions
% on team decision making, developmental economics, time preferences,
% and workplace interventions into the EBF framework.
%
% Language: English
% =============================================================================

% =============================================================================
% CROSS-REFERENCE STRUCTURE
% =============================================================================
%
% CHAPTER LINKAGE (Required):
% - Primary Chapter: Chapter 9: Complementarity and Interaction (CORE-HOW)
% - Secondary Chapters: Chapter 5 (WHO), Chapter 10 (Context), Chapter 14 (HR)
%
% This appendix provides evidence and parameters for:
%
% DEPENDENCIES (what this appendix REQUIRES):
% - Appendix UNMAPPED_B (CORE-HOW): Complementarity framework
% - Appendix UNMAPPED_AAA (CORE-WHO): Welfare hierarchy
% - Appendix UNMAPPED_V (CORE-WHEN): Context framework
% - Appendix K (LIT-FEHR): Social preferences foundation
%
% DEPENDENTS (what REQUIRES this appendix):
% - Appendix UNMAPPED_BBB (CORE-WHERE): Parameters for LLMMC
% - Appendix AH (DOMAIN-HR): Workplace interventions
% - All DOMAIN appendices using team-based interventions
%
% =============================================================================

\begin{tcolorbox}[colback=gray!5!white, colframe=gray!75!black,
    title=Appendix Metadata]
\textbf{Appendix:} SU --- LIT-SUTTER \\
\textbf{Category:} Literature Integration \\
\textbf{Author Focus:} Matthias Sutter (Max Planck Institute, Bonn) \\
\textbf{Papers:} 20 (90\% Tier-1 Evidence) \\
\textbf{Primary Contribution:} Team decision making, developmental economics
\end{tcolorbox}

\begin{tcolorbox}[colback=purple!5!white, colframe=purple!75!black,
    title=Chapter Linkage]

\textbf{Primary Chapter:} Chapter 9: Complementarity and Interaction
\begin{itemize}[nosep]
\item Section 9.3: Social Interaction $\leftrightarrow$ This Appendix Section \ref{sec:su-team}
\item Section 9.5: Group vs Individual $\leftrightarrow$ This Appendix Section \ref{sec:su-groupind}
\end{itemize}

\textbf{Secondary Chapters:}
\begin{itemize}[nosep]
\item Chapter 5: Welfare Hierarchy --- developmental trajectory of preferences
\item Chapter 10: Context --- institutional context $\Psi_I$
\item Chapter 14: HR Applications --- workplace interventions
\end{itemize}

\textbf{Reading Guide:}
\begin{itemize}[nosep]
\item For conceptual overview: Read Chapter 9 first
\item For parameter values: See Section \ref{sec:su-parameters}
\item For practical applications: See Section \ref{sec:su-practice}
\end{itemize}

\end{tcolorbox}


\chapter{LIT-SUTTER: Team Decision Making and Developmental Economics}
\label{app:lit-sutter}

\begin{abstract}
This appendix integrates Matthias Sutter's research into the EBF framework, covering four major themes: (1) team versus individual decision making, (2) developmental trajectory of economic preferences, (3) intergenerational transmission of preferences, and (4) scalable workplace interventions. Sutter's work provides critical evidence for CORE-HOW (complementarity in groups), CORE-WHO (age as welfare dimension), and CORE-WHEN (institutional context). With 90\% Tier-1 evidence quality (18 of 20 papers), this literature offers some of the most robust parameters for LLMMC estimation.
\end{abstract}


\begin{tcolorbox}[colback=blue!5!white, colframe=blue!75!black,
    title=Quick Reference: Matthias Sutter's Contributions]

\textbf{Core Question:} How do groups differ from individuals in economic decision making, and how do preferences develop across the lifespan?

\textbf{Key Finding:}
\begin{equation}
\gamma_{\text{group}} \neq \gamma_{\text{individual}} \quad \text{(Groups are less ``behavioral'')}
\end{equation}

\textbf{Key Concepts:}
\begin{itemize}[nosep]
\item \textbf{Team Rationality}: Groups converge faster to Nash equilibrium
\item \textbf{Developmental Trajectory}: Trust increases linearly from childhood
\item \textbf{Conditional Cooperation}: $\approx 50\%$ universal across cultures
\item \textbf{Workplace Climate}: Scalable interventions ($d = 0.25$)
\end{itemize}

\textbf{Cross-References:} Appendix UNMAPPED_B (CORE-HOW), K (LIT-FEHR), AAA (CORE-WHO), V (CORE-WHEN)
\end{tcolorbox}

% -----------------------------------------------------------------------------
% APPENDIX SCOPE BOX
% -----------------------------------------------------------------------------

\begin{tcolorbox}[
    colback=orange!5!white,
    colframe=orange!75!black,
    title={\textbf{Appendix Scope: Ziel / In-Scope / Out-of-Scope / Constraints}},
    fonttitle=\bfseries
]

\textbf{Ziel (Objective):}
\begin{itemize}[nosep]
    \item Do behavioral parameters transfer from individuals to groups, and how do preferences develop across the lifespan?
\end{itemize}

\textbf{In-Scope:}
\begin{itemize}[nosep]
    \item The Fundamental Question
    \item Theoretical Framework
    \item Axioms and Formal Definitions
    \item Key Results and Findings
    \item Theme 1: Team vs Individual Decision Making
\end{itemize}

\textbf{Out-of-Scope (delegiert an andere Appendices):}
\begin{itemize}[nosep]
    \item CORE-HOW $\rightarrow$ Appendix UNMAPPED_B
    \item CORE-WHO $\rightarrow$ Appendix UNMAPPED_AAA
    \item CORE-WHEN $\rightarrow$ Appendix UNMAPPED_V
    \item LIT-FEHR $\rightarrow$ Appendix K
    \item CORE-WHERE $\rightarrow$ Appendix UNMAPPED_BBB
\end{itemize}

\textbf{Constraints (Anwendungsgrenzen):}
\begin{itemize}[nosep]
    \item [TODO: Define when this appendix does NOT apply]
    \item [TODO: Define required prerequisites]
    \item [TODO: Define known limitations]
\end{itemize}

\textbf{Lieferobjekte (Deliverables):}
\begin{itemize}[nosep]
    \item 5 Table(s)
    \item 18 Equation(s)
    \item 6 Axiom(s)
    \item Worked Example(s)
\end{itemize}

\end{tcolorbox}




%%%%%%%%%%%%%%%%%%%%%%%%%%%%%%%%%%%%%%%%%%%%%%%%%%%%%%%%%%%%%%%%%%%%%%%%%%%%%%%
\section{The Fundamental Question}
\label{sec:su-fundamental}
%%%%%%%%%%%%%%%%%%%%%%%%%%%%%%%%%%%%%%%%%%%%%%%%%%%%%%%%%%%%%%%%%%%%%%%%%%%%%%%

\subsection{Why This Matters for EBF}

The EBF framework defines utility at multiple levels (individual $\ell = 0$, dyad $\ell = 1$, group $\ell = 2$, society $\ell = 3$). A critical question is whether behavioral parameters calibrated on individuals transfer to groups. Sutter's research program systematically addresses this gap.

Additionally, most behavioral economics research assumes adult preferences as given. Sutter's developmental work traces the \textit{emergence} of preferences from childhood, providing crucial evidence for CORE-WHO (who counts in welfare calculations) and heterogeneity in $\Psi_D$ (demographic context).

\subsection{The Core Problem}

Standard behavioral economics imports parameters from individual laboratory experiments into policy models. But:
\begin{itemize}
\item Most economic decisions involve groups (families, teams, committees)
\item Preferences develop over time---adults are not born with $\beta = 0.7$
\item Cross-cultural variation is often assumed but rarely tested
\end{itemize}

\begin{tcolorbox}[colback=red!5!white, colframe=red!75!black,
    title=The Fundamental Question]
\textbf{Do behavioral parameters transfer from individuals to groups, and how do preferences develop across the lifespan?}
\end{tcolorbox}


%%%%%%%%%%%%%%%%%%%%%%%%%%%%%%%%%%%%%%%%%%%%%%%%%%%%%%%%%%%%%%%%%%%%%%%%%%%%%%%
\section{Theoretical Framework}
\label{sec:su-theory}
%%%%%%%%%%%%%%%%%%%%%%%%%%%%%%%%%%%%%%%%%%%%%%%%%%%%%%%%%%%%%%%%%%%%%%%%%%%%%%%

Sutter's research program addresses three theoretical gaps in behavioral economics:

\subsection{The Group-Individual Gap}

Standard models assume behavioral parameters $(\beta, \gamma, \lambda)$ measured on individuals transfer directly to group settings. Sutter's systematic comparison shows this assumption fails:
\begin{equation}
\theta_{\text{group}} = f(\theta_{\text{individual}}, n, \text{deliberation})
\end{equation}
where $n$ is group size and deliberation captures the within-group discussion process.

\subsection{The Developmental Gap}

Most behavioral economics treats preferences as fixed adult characteristics. Sutter traces the \textit{ontogeny} of economic preferences:
\begin{equation}
\theta(a) = \theta_0 + g(a) + \epsilon \quad \text{(preference development with age } a \text{)}
\end{equation}

\subsection{The Transmission Gap}

Where do preferences come from? Sutter's intergenerational work shows family transmission:
\begin{equation}
\theta_{\text{child}} = \rho \cdot \theta_{\text{parent}} + (1-\rho) \cdot \theta_{\text{social}} + \epsilon
\end{equation}


%%%%%%%%%%%%%%%%%%%%%%%%%%%%%%%%%%%%%%%%%%%%%%%%%%%%%%%%%%%%%%%%%%%%%%%%%%%%%%%
\section{Axioms and Formal Definitions}
\label{sec:su-axioms}
%%%%%%%%%%%%%%%%%%%%%%%%%%%%%%%%%%%%%%%%%%%%%%%%%%%%%%%%%%%%%%%%%%%%%%%%%%%%%%%

Based on Sutter's empirical findings, we formalize six axioms for the EBF:

\begin{tcolorbox}[colback=gray!5!white, colframe=gray!75!black,
    title=Axiom UNMAPPED_SU-1: Group Rationality Premium]
\textbf{Statement:} In strategic interactions, groups exhibit higher levels of game-theoretic rationality than individuals.
\begin{equation}
\mathbb{E}[\pi_{\text{group}}^{\text{Nash}}] > \mathbb{E}[\pi_{\text{individual}}^{\text{Nash}}]
\end{equation}
\textbf{Interpretation:} Groups filter individual biases through deliberation.

\textbf{Source:} \citet{charness2012groups}, \citet{kocher2005decision}
\end{tcolorbox}

\begin{tcolorbox}[colback=gray!5!white, colframe=gray!75!black,
    title=Axiom UNMAPPED_SU-2: Group Selfishness Premium]
\textbf{Statement:} In allocation decisions, groups are more self-interested than individuals.
\begin{equation}
\gamma_{\text{group}} > \gamma_{\text{individual}} \quad \text{(less other-regarding)}
\end{equation}
\textbf{Interpretation:} Group deliberation amplifies self-interest while filtering altruistic biases.

\textbf{Source:} \citet{luhan2009polarization}, \citet{kugler2007trust}
\end{tcolorbox}

\begin{tcolorbox}[colback=gray!5!white, colframe=gray!75!black,
    title=Axiom UNMAPPED_SU-3: Trust Development]
\textbf{Statement:} Trust increases monotonically from childhood to early adulthood, then stabilizes.
\begin{equation}
\frac{\partial \text{Trust}}{\partial a} > 0 \text{ for } a < 20, \quad \frac{\partial \text{Trust}}{\partial a} \approx 0 \text{ for } a \geq 20
\end{equation}
\textbf{Interpretation:} Social capital accumulates through experience until maturity.

\textbf{Source:} \citet{sutter2007trust}
\end{tcolorbox}

\begin{tcolorbox}[colback=gray!5!white, colframe=gray!75!black,
    title=Axiom UNMAPPED_SU-4: Preference Heritability]
\textbf{Statement:} Economic preferences are partially transmitted within families.
\begin{equation}
\text{Corr}(\theta_{\text{parent}}, \theta_{\text{child}}) = \rho \in [0.3, 0.4]
\end{equation}
\textbf{Interpretation:} Nature and nurture both contribute to preference formation.

\textbf{Source:} \citet{chowdhury2022preferences}
\end{tcolorbox}

\begin{tcolorbox}[colback=gray!5!white, colframe=gray!75!black,
    title=Axiom UNMAPPED_SU-5: Conditional Cooperation Universality]
\textbf{Statement:} The share of conditional cooperators is approximately constant across cultures.
\begin{equation}
P(\text{CC}) \approx 0.50 \quad \forall \text{ cultures}
\end{equation}
\textbf{Interpretation:} Conditional cooperation is a human universal, not culturally contingent.

\textbf{Source:} \citet{kocher2008conditional}
\end{tcolorbox}

\begin{tcolorbox}[colback=gray!5!white, colframe=gray!75!black,
    title=Axiom UNMAPPED_SU-6: Lab-Field Validity]
\textbf{Statement:} Experimentally measured time preferences predict real-world behavior.
\begin{equation}
\text{Corr}(\beta_{\text{lab}}, \text{Behavior}_{\text{field}}) > 0.3
\end{equation}
\textbf{Interpretation:} Laboratory experiments have external validity for preference measurement.

\textbf{Source:} \citet{sutter2013impatience}
\end{tcolorbox}


%%%%%%%%%%%%%%%%%%%%%%%%%%%%%%%%%%%%%%%%%%%%%%%%%%%%%%%%%%%%%%%%%%%%%%%%%%%%%%%
\section{Key Results and Findings}
\label{sec:su-results}
%%%%%%%%%%%%%%%%%%%%%%%%%%%%%%%%%%%%%%%%%%%%%%%%%%%%%%%%%%%%%%%%%%%%%%%%%%%%%%%

This section consolidates the main empirical findings across all five themes:

\begin{table}[htbp]
\centering
\caption{Summary of Key Results from Sutter's Research}
\label{tab:su-key-results}
\renewcommand{\arraystretch}{1.3}
\begin{tabular}{@{}llcc@{}}
\toprule
\textbf{Result} & \textbf{Finding} & \textbf{Effect} & \textbf{Tier} \\
\midrule
UNMAPPED_SU-1 & Groups are less behavioral & $+0.30$ rationality & 1 \\
UNMAPPED_SU-2 & Groups trust less than individuals & $-0.20$ trust & 1 \\
UNMAPPED_SU-3 & Trust increases with age & $+0.15$ per stage & 1 \\
UNMAPPED_SU-4 & Time preferences predict health & $r = 0.35$ & 1 \\
UNMAPPED_SU-5 & Preferences transmitted in families & $\rho = 0.35$--$0.40$ & 1 \\
UNMAPPED_SU-6 & Conditional cooperation is universal & $\approx 50\%$ & 1 \\
UNMAPPED_SU-7 & Workplace training improves climate & $d = 0.25$ & 1 \\
UNMAPPED_SU-8 & Psychological pressure affects performance & $d = 0.21$ & 1 \\
\bottomrule
\end{tabular}
\end{table}


%%%%%%%%%%%%%%%%%%%%%%%%%%%%%%%%%%%%%%%%%%%%%%%%%%%%%%%%%%%%%%%%%%%%%%%%%%%%%%%
\section{Theme 1: Team vs Individual Decision Making}
\label{sec:su-team}
%%%%%%%%%%%%%%%%%%%%%%%%%%%%%%%%%%%%%%%%%%%%%%%%%%%%%%%%%%%%%%%%%%%%%%%%%%%%%%%

\subsection{The Central Finding}
\label{sec:su-groupind}

Sutter's research with Gary Charness \citep{charness2012groups} establishes a fundamental insight:

\begin{tcolorbox}[colback=green!5!white, colframe=green!75!black,
    title=Key Result UNMAPPED_SU-1: Groups are Less Behavioral]
\textbf{Statement:} Groups are more likely to follow game-theoretic predictions than individuals. Individuals are more influenced by biases, cognitive limitations, and social considerations.

\textbf{Implication for EBF:} The complementarity parameter $\gamma$ differs systematically between individual and group decision-making contexts:
\begin{equation}
\gamma_{\text{group}} > \gamma_{\text{individual}} \quad \text{(groups more self-interested)}
\end{equation}
\end{tcolorbox}

\subsection{Evidence from Beauty-Contest Games}

\citet{kocher2005decision} show that in experimental beauty-contest games:
\begin{itemize}
\item Groups converge faster to Nash equilibrium than individuals
\item Effect size: $d = 0.30$ improvement in strategic sophistication
\item Mechanism: Within-group deliberation filters out irrational choices
\end{itemize}

\citet{sutter2005four} extends this finding:
\begin{itemize}
\item Four-person teams are \textit{not} better than two-person teams
\item Diminishing returns to team size in strategic tasks
\item Optimal team size depends on task complexity
\end{itemize}

\subsection{Trust Between Groups}

\citet{kugler2007trust} finds an asymmetric pattern:
\begin{itemize}
\item Groups trust \textit{less} than individuals (effect: $-0.20$)
\item Groups are \textit{equally trustworthy} as individuals (effect: $+0.02$)
\end{itemize}

This has important implications for multi-level utility:
\begin{equation}
\text{Trust}_{\ell=2} < \text{Trust}_{\ell=0} \quad \text{but} \quad \text{Trustworthiness}_{\ell=2} \approx \text{Trustworthiness}_{\ell=0}
\end{equation}

\subsection{Group Polarization}

\citet{luhan2009polarization} shows:
\begin{itemize}
\item Teams are more selfish than individuals in dictator games
\item The most selfish team member has the strongest influence
\item Effect size: $d = -0.25$ (more selfish)
\end{itemize}

\begin{tcolorbox}[colback=yellow!5!white, colframe=yellow!75!black,
    title=Implication for LLMMC]
When modeling group decisions, use:
\begin{equation}
\gamma_{\text{group}} = \gamma_{\text{individual}} + 0.15 \quad \text{(toward self-interest)}
\end{equation}
This adjustment accounts for the ``wisdom of groups'' effect where groups filter biases but also amplify self-interest.
\end{tcolorbox}


%%%%%%%%%%%%%%%%%%%%%%%%%%%%%%%%%%%%%%%%%%%%%%%%%%%%%%%%%%%%%%%%%%%%%%%%%%%%%%%
\section{Theme 2: Developmental Trajectory of Preferences}
\label{sec:su-developmental}
%%%%%%%%%%%%%%%%%%%%%%%%%%%%%%%%%%%%%%%%%%%%%%%%%%%%%%%%%%%%%%%%%%%%%%%%%%%%%%%

\subsection{Trust Across Age Groups}

\citet{sutter2007trust} conducted trust games with 662 participants across six age groups (8 years to retirement):

\begin{table}[htbp]
\centering
\caption{Trust and Trustworthiness Across Age Groups}
\label{tab:su-trust-age}
\renewcommand{\arraystretch}{1.3}
\begin{tabular}{@{}lccc@{}}
\toprule
\textbf{Age Group} & \textbf{Trust (Send)} & \textbf{Trustworthiness (Return)} & \textbf{N} \\
\midrule
8 years & Low & Moderate & 110 \\
12 years & Medium-Low & Moderate & 110 \\
16 years & Medium & Moderate-High & 110 \\
Young adult (20s) & High & High & 112 \\
Middle-aged (30-50) & High & High & 110 \\
Retired (60+) & High & High & 110 \\
\bottomrule
\end{tabular}
\end{table}

\begin{tcolorbox}[colback=green!5!white, colframe=green!75!black,
    title=Key Result UNMAPPED_SU-2: Age-Trust Gradient]
Trust increases \textbf{linearly} from childhood to early adulthood, then plateaus. Trustworthiness is relatively stable across ages.
\begin{equation}
\text{Trust}(a) = \text{Trust}_0 + 0.15 \cdot \mathbb{1}_{a \geq 20} \quad (a = \text{age})
\end{equation}
\end{tcolorbox}

\subsection{Economic Behavior of Children: Survey Evidence}

\citet{sutter2019economic} provides a comprehensive survey showing:
\begin{itemize}
\item Risk preferences: Children are more risk-seeking; decreases with age
\item Time preferences: Children are more impatient; patience increases with age
\item Social preferences: Conditional cooperation emerges by age 6-7
\end{itemize}

\subsection{Impatience Predicts Field Behavior}

The landmark paper \citet{sutter2013impatience} (AER) demonstrates:

\begin{tcolorbox}[colback=blue!5!white, colframe=blue!75!black,
    title=Key Result UNMAPPED_SU-3: Lab-Field Validity]
Experimental measures of time preferences predict real-world behavior years later:
\begin{itemize}[nosep]
\item Impatient adolescents: higher BMI ($r = 0.35$)
\item Impatient adolescents: more likely to smoke/drink
\item Impatient adolescents: lower savings
\end{itemize}
\textbf{Implication:} Laboratory-measured $\beta$ has external validity.
\end{tcolorbox}


%%%%%%%%%%%%%%%%%%%%%%%%%%%%%%%%%%%%%%%%%%%%%%%%%%%%%%%%%%%%%%%%%%%%%%%%%%%%%%%
\section{Theme 3: Intergenerational Transmission}
\label{sec:su-intergenerational}
%%%%%%%%%%%%%%%%%%%%%%%%%%%%%%%%%%%%%%%%%%%%%%%%%%%%%%%%%%%%%%%%%%%%%%%%%%%%%%%

\subsection{Family Clustering of Preferences}

\citet{chowdhury2022preferences} (JPE) conducted large-scale experiments in Bangladesh (N = 1,999 individuals from 544 families):

\begin{tcolorbox}[colback=green!5!white, colframe=green!75!black,
    title=Key Result UNMAPPED_SU-4: Preference Heritability]
\textbf{Finding:} Strong intergenerational persistence of economic preferences
\begin{itemize}[nosep]
\item Correlation between parent and child time preferences: $r = 0.35$
\item Correlation between parent and child social preferences: $r = 0.40$
\item Families cluster into types: Patient/Prosocial vs Impatient/Selfish
\end{itemize}
\textbf{Implication for $\Psi_D$:} Family background is a critical context dimension.
\end{tcolorbox}

This finding has major implications for CORE-WHO:
\begin{equation}
\text{Preferences}_{t+1} = \rho \cdot \text{Preferences}_t^{\text{parents}} + \epsilon
\end{equation}
where $\rho \approx 0.35$--$0.40$.


%%%%%%%%%%%%%%%%%%%%%%%%%%%%%%%%%%%%%%%%%%%%%%%%%%%%%%%%%%%%%%%%%%%%%%%%%%%%%%%
\section{Theme 4: Cross-Cultural Universals}
\label{sec:su-crosscultural}
%%%%%%%%%%%%%%%%%%%%%%%%%%%%%%%%%%%%%%%%%%%%%%%%%%%%%%%%%%%%%%%%%%%%%%%%%%%%%%%

\subsection{Conditional Cooperation Across Continents}

\citet{kocher2008conditional} tested conditional cooperation in USA, Austria, and Japan:

\begin{tcolorbox}[colback=green!5!white, colframe=green!75!black,
    title=Key Result UNMAPPED_SU-5: Universal Conditional Cooperation]
\textbf{Finding:} Approximately 50\% of subjects are conditional cooperators across all three cultures.
\begin{equation}
P(\text{Conditional Cooperator}) \approx 0.50 \quad \text{(universal)}
\end{equation}
\textbf{Implication:} This parameter can be used as a robust default in LLMMC.
\end{tcolorbox}

\subsection{Language and Time Preferences}

\citet{sutter2018language} exploits natural variation in a bilingual city (South Tyrol):
\begin{itemize}
\item German-speaking children: more patient
\item Italian-speaking children: more impatient
\item Effect size: $d = 0.18$
\end{itemize}

This provides evidence for $\Psi_C$ (cultural context) affecting preferences.


%%%%%%%%%%%%%%%%%%%%%%%%%%%%%%%%%%%%%%%%%%%%%%%%%%%%%%%%%%%%%%%%%%%%%%%%%%%%%%%
\section{Theme 5: Workplace Interventions}
\label{sec:su-workplace}
%%%%%%%%%%%%%%%%%%%%%%%%%%%%%%%%%%%%%%%%%%%%%%%%%%%%%%%%%%%%%%%%%%%%%%%%%%%%%%%

\subsection{Improving Workplace Climate at Scale}

\citet{alan2023workplace} (QJE) is one of the largest workplace RCTs:
\begin{itemize}
\item Setting: 24 large corporations in Turkey
\item N: 2,000+ white-collar professionals
\item Design: Clustered randomized intervention
\end{itemize}

Results:
\begin{itemize}
\item Workplace climate improvement: $d = 0.25$
\item Turnover reduction: significant
\item Social networks: denser and less segregated ($d = 0.30$)
\end{itemize}

\begin{tcolorbox}[colback=green!5!white, colframe=green!75!black,
    title=Key Result UNMAPPED_SU-6: Scalable Workplace Interventions]
Training-based interventions can improve workplace climate in large organizations.
\begin{equation}
\Delta \text{Climate} = 0.25 \cdot \text{Treatment} \quad (p < 0.01)
\end{equation}
\end{tcolorbox}

\subsection{Female Leadership Effects}

\citet{alan2024female} shows:
\begin{itemize}
\item Female leaders improve workplace climate for female employees
\item Social support from leaders moderates perceptions
\item Gender-congruent leadership has positive effects
\end{itemize}

\subsection{Adverse Employer Behavior}

\citet{heinz2020indirect} demonstrates third-party fairness effects:
\begin{itemize}
\item Workers reduce effort when observing unfair treatment of coworkers
\item Effect size: $d = -0.15$
\item Mechanism: Fairness concerns extend beyond direct recipients
\end{itemize}


%%%%%%%%%%%%%%%%%%%%%%%%%%%%%%%%%%%%%%%%%%%%%%%%%%%%%%%%%%%%%%%%%%%%%%%%%%%%%%%
\section{Parameters for LLMMC}
\label{sec:su-parameters}
%%%%%%%%%%%%%%%%%%%%%%%%%%%%%%%%%%%%%%%%%%%%%%%%%%%%%%%%%%%%%%%%%%%%%%%%%%%%%%%

\begin{table}[htbp]
\centering
\caption{Sutter Parameters for LLMMC Estimation}
\label{tab:su-parameters}
\renewcommand{\arraystretch}{1.3}
\begin{tabular}{@{}llcc@{}}
\toprule
\textbf{Parameter} & \textbf{Value} & \textbf{Tier} & \textbf{Source} \\
\midrule
\multicolumn{4}{l}{\textit{Team Decision Making}} \\
Group rationality improvement & $+0.30$ & 1 & kocher2005decision \\
Group trust deficit & $-0.20$ & 1 & kugler2007trust \\
Group selfishness increase & $-0.25$ & 1 & luhan2009polarization \\
\midrule
\multicolumn{4}{l}{\textit{Developmental}} \\
Impatience-health correlation & $0.35$ & 1 & sutter2013impatience \\
Trust age gradient & $+0.15$ & 1 & sutter2007trust \\
\midrule
\multicolumn{4}{l}{\textit{Intergenerational}} \\
Preference heritability & $0.35$--$0.40$ & 1 & chowdhury2022preferences \\
\midrule
\multicolumn{4}{l}{\textit{Cross-Cultural}} \\
Conditional cooperation rate & $0.50$ & 1 & kocher2008conditional \\
Language-patience effect & $0.18$ & 1 & sutter2018language \\
\midrule
\multicolumn{4}{l}{\textit{Workplace}} \\
Climate improvement (training) & $0.25$ & 1 & alan2023workplace \\
Adverse behavior spillover & $-0.15$ & 1 & heinz2020indirect \\
Pressure-performance effect & $0.21$ & 1 & kocher2012pressure \\
\bottomrule
\end{tabular}
\end{table}


%%%%%%%%%%%%%%%%%%%%%%%%%%%%%%%%%%%%%%%%%%%%%%%%%%%%%%%%%%%%%%%%%%%%%%%%%%%%%%%
\section{Integration with EBF}
\label{sec:su-integration}
%%%%%%%%%%%%%%%%%%%%%%%%%%%%%%%%%%%%%%%%%%%%%%%%%%%%%%%%%%%%%%%%%%%%%%%%%%%%%%%

\begin{tcolorbox}[colback=yellow!5!white, colframe=yellow!75!black,
    title=Integration: LIT-SUTTER $\leftrightarrow$ CORE Appendices]

\textbf{Connection to CORE-HOW (Appendix UNMAPPED_B):}
\begin{itemize}[nosep]
\item Provides evidence that $\gamma_{\text{group}} \neq \gamma_{\text{individual}}$
\item Quantifies the ``groups are less behavioral'' effect
\item Informs multi-level utility specifications
\end{itemize}

\textbf{Connection to CORE-WHO (Appendix UNMAPPED_AAA):}
\begin{itemize}[nosep]
\item Establishes developmental trajectory of preferences
\item Justifies age as critical welfare dimension
\item Documents intergenerational transmission
\end{itemize}

\textbf{Connection to CORE-WHEN (Appendix UNMAPPED_V):}
\begin{itemize}[nosep]
\item $\Psi_I$ (Institutional): Team vs individual contexts
\item $\Psi_D$ (Demographic): Age and family background
\item $\Psi_C$ (Cultural): Language effects on preferences
\item $\Psi_P$ (Physical): Pressure and competition effects
\end{itemize}

\textbf{Connection to CORE-WHERE (Appendix UNMAPPED_BBB):}
\begin{itemize}[nosep]
\item 11 validated parameters for LLMMC
\item 90\% Tier-1 evidence quality
\item Cross-cultural replication
\end{itemize}

\end{tcolorbox}


%%%%%%%%%%%%%%%%%%%%%%%%%%%%%%%%%%%%%%%%%%%%%%%%%%%%%%%%%%%%%%%%%%%%%%%%%%%%%%%
\section{Worked Example}
\label{sec:su-example}
%%%%%%%%%%%%%%%%%%%%%%%%%%%%%%%%%%%%%%%%%%%%%%%%%%%%%%%%%%%%%%%%%%%%%%%%%%%%%%%

\subsection{Designing a Team-Based Savings Intervention}

\paragraph{Setup.} A company wants to increase retirement savings among employees. Should the intervention target individuals or teams?

\paragraph{Application.} Using Sutter's findings:
\begin{enumerate}
\item Individual intervention: Standard nudges (defaults, reminders)
\item Team intervention: Group discussions about retirement planning
\end{enumerate}

\paragraph{Expected Effects:}
\begin{itemize}
\item Teams will be more ``rational'' (less influenced by present bias)
\item But teams will also be more self-interested (may resist employer-favored options)
\item Net effect depends on whether present bias or strategic selfishness dominates
\end{itemize}

\paragraph{Recommendation:}
For retirement savings, where present bias is the main obstacle:
\begin{equation}
\text{Team-based intervention} \succ \text{Individual nudge}
\end{equation}
because teams filter present-bias irrationality (Sutter's main finding).


%%%%%%%%%%%%%%%%%%%%%%%%%%%%%%%%%%%%%%%%%%%%%%%%%%%%%%%%%%%%%%%%%%%%%%%%%%%%%%%
\section{Implications for Practice}
\label{sec:su-practice}
%%%%%%%%%%%%%%%%%%%%%%%%%%%%%%%%%%%%%%%%%%%%%%%%%%%%%%%%%%%%%%%%%%%%%%%%%%%%%%%

\begin{tcolorbox}[colback=green!5!white, colframe=green!75!black,
    title=Practical Implications]
\begin{enumerate}
\item \textbf{Team-Based Interventions}: When the goal is to overcome biases, group deliberation can help. But expect increased self-interest.

\item \textbf{Age-Appropriate Design}: Interventions for children need to account for lower trust and higher impatience. Don't use adult parameters.

\item \textbf{Family-Level Targeting}: Since preferences cluster in families, targeting family units may be more efficient than targeting individuals.

\item \textbf{Cross-Cultural Robustness}: The 50\% conditional cooperator finding is universal---use this as default when culture-specific data is unavailable.

\item \textbf{Workplace Climate}: Training interventions can improve climate at scale ($d = 0.25$). Consider female leadership for female employees.
\end{enumerate}
\end{tcolorbox}


%%%%%%%%%%%%%%%%%%%%%%%%%%%%%%%%%%%%%%%%%%%%%%%%%%%%%%%%%%%%%%%%%%%%%%%%%%%%%%%
\section{Summary}
\label{sec:su-summary}
%%%%%%%%%%%%%%%%%%%%%%%%%%%%%%%%%%%%%%%%%%%%%%%%%%%%%%%%%%%%%%%%%%%%%%%%%%%%%%%

\begin{tcolorbox}[colback=green!10!white, colframe=green!75!black,
    title=\textbf{Appendix UNMAPPED_SU: Summary}]

\textbf{Core Question:} How do groups differ from individuals, and how do preferences develop?

\textbf{Main Contributions:}
\begin{itemize}[nosep]
    \item Groups are less ``behavioral'' but more self-interested
    \item Trust develops linearly from childhood to adulthood
    \item Preferences are transmitted within families ($r = 0.35$--$0.40$)
    \item Conditional cooperation ($\approx 50\%$) is universal
    \item Workplace climate can be improved at scale
\end{itemize}

\textbf{Key Parameters:}
\begin{itemize}[nosep]
    \item $\gamma_{\text{group}} = \gamma_{\text{individual}} + 0.15$
    \item Trust age gradient: $+0.15$
    \item Conditional cooperation: $0.50$
    \item Workplace intervention: $d = 0.25$
\end{itemize}

\textbf{Evidence Quality:} 90\% Tier-1 (highest in EBF literature appendices)

\textbf{Integration:} Connects to CORE-HOW, CORE-WHO, CORE-WHEN, CORE-WHERE

\end{tcolorbox}


%%%%%%%%%%%%%%%%%%%%%%%%%%%%%%%%%%%%%%%%%%%%%%%%%%%%%%%%%%%%%%%%%%%%%%%%%%%%%%%
\section{Glossary of Symbols}
\label{sec:su-glossary}
%%%%%%%%%%%%%%%%%%%%%%%%%%%%%%%%%%%%%%%%%%%%%%%%%%%%%%%%%%%%%%%%%%%%%%%%%%%%%%%

\begin{tcolorbox}[colback=blue!5!white, colframe=blue!75!black]
\textbf{Central Reference:} For complete symbol definitions, see
\textbf{Appendix G} (\textit{Glossary and Notation}).

This section lists only symbols introduced or specialized in this appendix.
\end{tcolorbox}

\begin{table}[htbp]
\centering
\caption{Symbols Introduced in Appendix UNMAPPED_SU}
\label{tab:su-symbols}
\renewcommand{\arraystretch}{1.3}
\begin{tabular}{@{}clp{6cm}c@{}}
\toprule
\textbf{Symbol} & \textbf{Name} & \textbf{Definition} & \textbf{Also in G?} \\
\midrule
$\gamma_{\text{group}}$ & Group complementarity & Complementarity parameter for group decisions & NEW \\
$\gamma_{\text{individual}}$ & Individual complementarity & Complementarity parameter for individual decisions & \checkmark \\
$\rho$ & Preference heritability & Parent-child preference correlation & NEW \\
\bottomrule
\end{tabular}
\end{table}


%%%%%%%%%%%%%%%%%%%%%%%%%%%%%%%%%%%%%%%%%%%%%%%%%%%%%%%%%%%%%%%%%%%%%%%%%%%%%%%
\section{Critical Foundations and Objections}
\label{sec:su-foundations}
%%%%%%%%%%%%%%%%%%%%%%%%%%%%%%%%%%%%%%%%%%%%%%%%%%%%%%%%%%%%%%%%%%%%%%%%%%%%%%%

\subsection{Objection 1: Laboratory Artificiality}

\textbf{Objection:} Laboratory experiments with students don't generalize to real-world group decisions in organizations.

\textbf{Response:} Sutter's research program systematically addresses this through:
\begin{itemize}
\item Field experiments in actual workplaces (QJE 2023: 24 corporations)
\item Lab-field linkages showing experimental measures predict real behavior (AER 2013)
\item Cross-cultural replications (USA, Austria, Japan, Bangladesh, Turkey)
\end{itemize}

\subsection{Objection 2: WEIRD Sample Bias}

\textbf{Objection:} Most experiments use Western, Educated, Industrialized, Rich, Democratic samples.

\textbf{Response:} Sutter explicitly tests cross-cultural validity:
\begin{itemize}
\item Conditional cooperation: replicated in USA, Austria, Japan \citep{kocher2008conditional}
\item Preference transmission: tested in Bangladesh \citep{chowdhury2022preferences}
\item Workplace interventions: tested in Turkey \citep{alan2023workplace}
\item The 50\% conditional cooperation finding is remarkably stable
\end{itemize}

\subsection{Objection 3: Selection vs Treatment Effects}

\textbf{Objection:} Group differences may reflect self-selection rather than causal effects of group deliberation.

\textbf{Response:} Sutter's designs address this through:
\begin{itemize}
\item Random assignment to individual vs group conditions
\item Within-subject comparisons \citep{luhan2009polarization}
\item Clustered randomization in field settings \citep{alan2023workplace}
\end{itemize}


%%%%%%%%%%%%%%%%%%%%%%%%%%%%%%%%%%%%%%%%%%%%%%%%%%%%%%%%%%%%%%%%%%%%%%%%%%%%%%%
\section{Open Issues and Research Agenda}
\label{sec:su-open}
%%%%%%%%%%%%%%%%%%%%%%%%%%%%%%%%%%%%%%%%%%%%%%%%%%%%%%%%%%%%%%%%%%%%%%%%%%%%%%%

\begin{table}[htbp]
\centering
\caption{Research Roadmap for Sutter-Related Questions}
\label{tab:su-roadmap}
\renewcommand{\arraystretch}{1.3}
\begin{tabular}{@{}clcl@{}}
\toprule
\textbf{\#} & \textbf{Issue} & \textbf{Priority} & \textbf{Status} \\
\midrule
1 & Optimal team size for different tasks & HIGH & Open \\
2 & Digital vs in-person team effects & HIGH & Open \\
3 & Long-run persistence of workplace interventions & MEDIUM & Partial \\
4 & Cross-cultural variation in group effects & MEDIUM & Partial \\
5 & Developmental trajectory beyond Western samples & LOW & Open \\
\bottomrule
\end{tabular}
\end{table}


%%%%%%%%%%%%%%%%%%%%%%%%%%%%%%%%%%%%%%%%%%%%%%%%%%%%%%%%%%%%%%%%%%%%%%%%%%%%%%%
\section{References}
\label{sec:su-references}
%%%%%%%%%%%%%%%%%%%%%%%%%%%%%%%%%%%%%%%%%%%%%%%%%%%%%%%%%%%%%%%%%%%%%%%%%%%%%%%

\begin{tcolorbox}[colback=blue!5!white, colframe=blue!75!black]
\textbf{Master Bibliography:} All references appear in \texttt{bcm\_master.bib}.

\textbf{Paper Database:} Full metadata in \texttt{data/paper-sources.yaml} with \texttt{lit\_appendix: LIT-SUTTER}.
\end{tcolorbox}

\subsection{Core References (20 Papers)}

\textbf{Team Decision Making:}
\begin{itemize}[nosep]
\item Charness \& Sutter (2012), JEP --- Groups are less behavioral
\item Kocher \& Sutter (2005), EJ --- Beauty-contest games
\item Sutter (2005), EconLett --- Team size effects
\item Kugler et al. (2007), JEP --- Trust between groups
\item Luhan et al. (2009), ExpEcon --- Group polarization
\item Sutter \& Strassmair (2009), GEB --- Communication effects
\end{itemize}

\textbf{Developmental Economics:}
\begin{itemize}[nosep]
\item Sutter et al. (2013), AER --- Impatience predicts field behavior
\item Sutter \& Kocher (2007), GEB --- Trust across ages
\item Sutter et al. (2019), EER --- Survey of children's behavior
\item Angerer et al. (2018), EER --- Language and patience
\end{itemize}

\textbf{Intergenerational:}
\begin{itemize}[nosep]
\item Chowdhury et al. (2022), JPE --- Family preference clusters
\end{itemize}

\textbf{Cross-Cultural:}
\begin{itemize}[nosep]
\item Kocher et al. (2008), EconLett --- Conditional cooperation worldwide
\end{itemize}

\textbf{Workplace:}
\begin{itemize}[nosep]
\item Alan et al. (2023), QJE --- Workplace climate RCT
\item Alan et al. (2024), ManSci --- Female leadership
\item Heinz et al. (2020), EJ --- Adverse employer behavior
\item Kocher et al. (2012), ManSci --- Psychological pressure
\end{itemize}

\textbf{Leadership:}
\begin{itemize}[nosep]
\item G\"uth et al. (2007), JPubE --- Leading by example
\item Rivas \& Sutter (2011), EconLett --- Voluntary leadership
\end{itemize}

\textbf{Social Preferences:}
\begin{itemize}[nosep]
\item Kirchler et al. (2023), ManSci --- Finance professionals
\item Barron et al. (2023), REStat --- Discrimination and narratives
\end{itemize}

\nocite{bcm_master}


%%%%%%%%%%%%%%%%%%%%%%%%%%%%%%%%%%%%%%%%%%%%%%%%%%%%%%%%%%%%%%%%%%%%%%%%%%%%%%%
% CROSS-REFERENCE FOOTER
%%%%%%%%%%%%%%%%%%%%%%%%%%%%%%%%%%%%%%%%%%%%%%%%%%%%%%%%%%%%%%%%%%%%%%%%%%%%%%%

\vspace{1em}
\hrule
\vspace{0.5em}

\noindent\textit{Cross-references:}
Appendix G (Central Glossary),
Appendix UNMAPPED_B (CORE-HOW),
Appendix UNMAPPED_AAA (CORE-WHO),
Appendix UNMAPPED_V (CORE-WHEN),
Appendix UNMAPPED_BBB (CORE-WHERE),
Appendix K (LIT-FEHR).

\noindent\textit{Master Bibliography:} \texttt{bcm\_master.bib}

% =============================================================================
% END OF APPENDIX SU: LIT-SUTTER
% =============================================================================
